\section{更新報酬定理の適用条件}

\begin{lemma}[潜伏状態の平均滞在時間の有限性]
$t_m>0,\ \beta(I)>0,\ \alpha(I)>0$ の任意の組に対し，
\[
  \E[\min(T_{\mathrm{det}},T_d)]
  =\frac{1-e^{-\beta(I)t_m}}{\beta(I)}
  +t_m^{\alpha(I)}\beta(I)^{\alpha(I)-1}
  \Gamma(1-\alpha(I),\beta(I)t_m)
  <\infty .
  \label{eq:min-finite}
\]
\end{lemma}

\begin{proof}
$\E[\min(X,Y)]=\int_0^\infty\Prob(X>t,Y>t)\dd t$ と独立性を用いると，
\[
  \E[\min(T_{\mathrm{det}},T_d)]
  =\int_0^{t_m}e^{-\beta t}\dd t
  +\int_{t_m}^{\infty}e^{-\beta t}\left(\frac{t_m}{t}\right)^\alpha\dd t .
\]
第一項は $(1-e^{-\beta t_m})/\beta$ である．第二項で $u=\beta t$ と置換すると
$t_m^\alpha\beta^{\alpha-1}\Gamma(1-\alpha,\beta t_m)$ を得る．
$x_0=\beta t_m>0$ に対し，$u\ge x_0$ では $u^{-\alpha}\le x_0^{-\alpha}$ なので，
\[
  \Gamma(1-\alpha,x_0)
  =\int_{x_0}^{\infty}u^{-\alpha}e^{-u}\dd u
  \le x_0^{-\alpha}e^{-x_0}<\infty .
\]
従って右辺は有限である．
\end{proof}

\begin{theorem}[更新報酬定理の適用可能性]
仮定 2.2 と 2.4 のもとで，任意の
$\lambda>0,\ \beta(I)>0,\ n\ge1,\ \Delta\in\mathcal{D}_n$ に対して，
合成半マルコフ過程は更新サイクル表現を持ち，サイクル長の期待値と
サイクルあたり期待費用は有限である．従って，更新報酬定理を適用でき，
\eqref{eq:cost-rate} は長期平均費用レートと一致する．
\end{theorem}

\begin{proof}
攻撃到着までの待ち時間は平均 $1/\lambda$ の指数分布である．
潜伏状態での滞在時間は補題 3.1 により
$\E[\min(T_{\mathrm{det}},T_d)]<\infty$ である．発症後の暗号化時間，
隔離，クローン，再デプロイの各時間は仮定により指数分布または
それらの有限和である．全損分岐と復旧完了分岐はいずれも費用を
計上した後に $S_N$ へ戻る更新時点として扱うため，サイクル長の
期待値は有限である．

費用については，全損費用 $L$，固定復旧費 $d_0$，期待復旧時間
$\tau_R$ は有限である．また $\Delta\in\mathcal{D}_n$ なら
$\Mbeta(\Delta;\alpha)$ は有限区間 $[t_m,n\Delta]$ 上の
切断積分であり有限である．早期検知分岐の費用 $\Qearly(I)$ は
$\Delta$ に依存しない有限な封じ込め費用として分離している．
従って一サイクルあたり期待費用は有限であり，更新報酬定理
\cite{cox1962,karlin1975,ross2014} を適用できる．
\end{proof}
